\IfFileExists{stacks-project.cls}{%
\documentclass{stacks-project}
}{%
\documentclass{amsart}
}

% For dealing with references we use the comment environment
\usepackage{verbatim}
\newenvironment{reference}{\comment}{\endcomment}
%\newenvironment{reference}{}{}
\newenvironment{slogan}{\comment}{\endcomment}
\newenvironment{history}{\comment}{\endcomment}

% For commutative diagrams we use Xy-pic
\usepackage[all]{xy}

% We use 2cell for 2-commutative diagrams.
\xyoption{2cell}
\UseAllTwocells

% We use multicol for the list of chapters between chapters
\usepackage{multicol}

% This is generally recommended for better output
\usepackage{lmodern}
\usepackage[T1]{fontenc}

% For cross-file-references
\usepackage{xr-hyper}

% Package for hypertext links:
\usepackage{hyperref}

% For any local file, say "hello.tex" you want to link to please
% use \externaldocument[hello-]{hello}
\externaldocument[introduction-]{introduction}
\externaldocument[conventions-]{conventions}
\externaldocument[sets-]{sets}
\externaldocument[categories-]{categories}
\externaldocument[topology-]{topology}
\externaldocument[sheaves-]{sheaves}
\externaldocument[sites-]{sites}
\externaldocument[stacks-]{stacks}
\externaldocument[fields-]{fields}
\externaldocument[algebra-]{algebra}
\externaldocument[brauer-]{brauer}
\externaldocument[homology-]{homology}
\externaldocument[derived-]{derived}
\externaldocument[simplicial-]{simplicial}
\externaldocument[more-algebra-]{more-algebra}
\externaldocument[smoothing-]{smoothing}
\externaldocument[modules-]{modules}
\externaldocument[sites-modules-]{sites-modules}
\externaldocument[injectives-]{injectives}
\externaldocument[cohomology-]{cohomology}
\externaldocument[sites-cohomology-]{sites-cohomology}
\externaldocument[dga-]{dga}
\externaldocument[dpa-]{dpa}
\externaldocument[sdga-]{sdga}
\externaldocument[hypercovering-]{hypercovering}
\externaldocument[schemes-]{schemes}
\externaldocument[constructions-]{constructions}
\externaldocument[properties-]{properties}
\externaldocument[morphisms-]{morphisms}
\externaldocument[coherent-]{coherent}
\externaldocument[divisors-]{divisors}
\externaldocument[limits-]{limits}
\externaldocument[varieties-]{varieties}
\externaldocument[topologies-]{topologies}
\externaldocument[descent-]{descent}
\externaldocument[perfect-]{perfect}
\externaldocument[more-morphisms-]{more-morphisms}
\externaldocument[flat-]{flat}
\externaldocument[groupoids-]{groupoids}
\externaldocument[more-groupoids-]{more-groupoids}
\externaldocument[etale-]{etale}
\externaldocument[chow-]{chow}
\externaldocument[intersection-]{intersection}
\externaldocument[pic-]{pic}
\externaldocument[weil-]{weil}
\externaldocument[adequate-]{adequate}
\externaldocument[dualizing-]{dualizing}
\externaldocument[duality-]{duality}
\externaldocument[discriminant-]{discriminant}
\externaldocument[derham-]{derham}
\externaldocument[local-cohomology-]{local-cohomology}
\externaldocument[algebraization-]{algebraization}
\externaldocument[curves-]{curves}
\externaldocument[resolve-]{resolve}
\externaldocument[models-]{models}
\externaldocument[functors-]{functors}
\externaldocument[equiv-]{equiv}
\externaldocument[pione-]{pione}
\externaldocument[etale-cohomology-]{etale-cohomology}
\externaldocument[proetale-]{proetale}
\externaldocument[relative-cycles-]{relative-cycles}
\externaldocument[more-etale-]{more-etale}
\externaldocument[trace-]{trace}
\externaldocument[crystalline-]{crystalline}
\externaldocument[spaces-]{spaces}
\externaldocument[spaces-properties-]{spaces-properties}
\externaldocument[spaces-morphisms-]{spaces-morphisms}
\externaldocument[decent-spaces-]{decent-spaces}
\externaldocument[spaces-cohomology-]{spaces-cohomology}
\externaldocument[spaces-limits-]{spaces-limits}
\externaldocument[spaces-divisors-]{spaces-divisors}
\externaldocument[spaces-over-fields-]{spaces-over-fields}
\externaldocument[spaces-topologies-]{spaces-topologies}
\externaldocument[spaces-descent-]{spaces-descent}
\externaldocument[spaces-perfect-]{spaces-perfect}
\externaldocument[spaces-more-morphisms-]{spaces-more-morphisms}
\externaldocument[spaces-flat-]{spaces-flat}
\externaldocument[spaces-groupoids-]{spaces-groupoids}
\externaldocument[spaces-more-groupoids-]{spaces-more-groupoids}
\externaldocument[bootstrap-]{bootstrap}
\externaldocument[spaces-pushouts-]{spaces-pushouts}
\externaldocument[spaces-chow-]{spaces-chow}
\externaldocument[groupoids-quotients-]{groupoids-quotients}
\externaldocument[spaces-more-cohomology-]{spaces-more-cohomology}
\externaldocument[spaces-simplicial-]{spaces-simplicial}
\externaldocument[spaces-duality-]{spaces-duality}
\externaldocument[formal-spaces-]{formal-spaces}
\externaldocument[restricted-]{restricted}
\externaldocument[spaces-resolve-]{spaces-resolve}
\externaldocument[formal-defos-]{formal-defos}
\externaldocument[defos-]{defos}
\externaldocument[cotangent-]{cotangent}
\externaldocument[examples-defos-]{examples-defos}
\externaldocument[algebraic-]{algebraic}
\externaldocument[examples-stacks-]{examples-stacks}
\externaldocument[stacks-sheaves-]{stacks-sheaves}
\externaldocument[criteria-]{criteria}
\externaldocument[artin-]{artin}
\externaldocument[quot-]{quot}
\externaldocument[stacks-properties-]{stacks-properties}
\externaldocument[stacks-morphisms-]{stacks-morphisms}
\externaldocument[stacks-limits-]{stacks-limits}
\externaldocument[stacks-cohomology-]{stacks-cohomology}
\externaldocument[stacks-perfect-]{stacks-perfect}
\externaldocument[stacks-introduction-]{stacks-introduction}
\externaldocument[stacks-more-morphisms-]{stacks-more-morphisms}
\externaldocument[stacks-geometry-]{stacks-geometry}
\externaldocument[moduli-]{moduli}
\externaldocument[moduli-curves-]{moduli-curves}
\externaldocument[examples-]{examples}
\externaldocument[exercises-]{exercises}
\externaldocument[guide-]{guide}
\externaldocument[desirables-]{desirables}
\externaldocument[coding-]{coding}
\externaldocument[obsolete-]{obsolete}
\externaldocument[fdl-]{fdl}
\externaldocument[index-]{index}

% Theorem environments.
%
\theoremstyle{plain}
\newtheorem{theorem}[subsection]{Theorem}
\newtheorem{proposition}[subsection]{Proposition}
\newtheorem{lemma}[subsection]{Lemma}

\theoremstyle{definition}
\newtheorem{definition}[subsection]{Definition}
\newtheorem{example}[subsection]{Example}
\newtheorem{exercise}[subsection]{Exercise}
\newtheorem{situation}[subsection]{Situation}

\theoremstyle{remark}
\newtheorem{remark}[subsection]{Remark}
\newtheorem{remarks}[subsection]{Remarks}

\numberwithin{equation}{subsection}

% Macros
%
\def\lim{\mathop{\mathrm{lim}}\nolimits}
\def\colim{\mathop{\mathrm{colim}}\nolimits}
\def\Spec{\mathop{\mathrm{Spec}}}
\def\Hom{\mathop{\mathrm{Hom}}\nolimits}
\def\Ext{\mathop{\mathrm{Ext}}\nolimits}
\def\SheafHom{\mathop{\mathcal{H}\!\mathit{om}}\nolimits}
\def\SheafExt{\mathop{\mathcal{E}\!\mathit{xt}}\nolimits}
\def\Sch{\mathit{Sch}}
\def\Mor{\mathop{\mathrm{Mor}}\nolimits}
\def\Ob{\mathop{\mathrm{Ob}}\nolimits}
\def\Sh{\mathop{\mathit{Sh}}\nolimits}
\def\NL{\mathop{N\!L}\nolimits}
\def\CH{\mathop{\mathrm{CH}}\nolimits}
\def\proetale{{pro\text{-}\acute{e}tale}}
\def\etale{{\acute{e}tale}}
\def\QCoh{\mathit{QCoh}}
\def\Ker{\mathop{\mathrm{Ker}}}
\def\Im{\mathop{\mathrm{Im}}}
\def\Coker{\mathop{\mathrm{Coker}}}
\def\Coim{\mathop{\mathrm{Coim}}}

% Boxtimes
%
\DeclareMathSymbol{\boxtimes}{\mathbin}{AMSa}{"02}

%
% Macros for moduli stacks/spaces
%
\def\QCohstack{\mathcal{QC}\!\mathit{oh}}
\def\Cohstack{\mathcal{C}\!\mathit{oh}}
\def\Spacesstack{\mathcal{S}\!\mathit{paces}}
\def\Quotfunctor{\mathrm{Quot}}
\def\Hilbfunctor{\mathrm{Hilb}}
\def\Curvesstack{\mathcal{C}\!\mathit{urves}}
\def\Polarizedstack{\mathcal{P}\!\mathit{olarized}}
\def\Complexesstack{\mathcal{C}\!\mathit{omplexes}}
% \Pic is the operator that assigns to X its picard group, usage \Pic(X)
% \Picardstack_{X/B} denotes the Picard stack of X over B
% \Picardfunctor_{X/B} denotes the Picard functor of X over B
\def\Pic{\mathop{\mathrm{Pic}}\nolimits}
\def\Picardstack{\mathcal{P}\!\mathit{ic}}
\def\Picardfunctor{\mathrm{Pic}}
\def\Deformationcategory{\mathcal{D}\!\mathit{ef}}
\usepackage{cleveref}
\crefname{originalreference}{original reference}{original references}
\Crefname{originalreference}{Original reference}{Original references}
\makeatletter
\newlabel{algebraic-lemma-base-change-by-space-representable-by-space}{{0300}{source}{}{sourcepointer-1}{}}
\newlabel{algebraic-lemma-base-change-by-space-representable-by-space@cref}{{[originalreference][0][]0300}{[1][0][]source}{}{}{}}
\makeatother
\crefname{originalreference}{original reference}{original references}
\Crefname{originalreference}{Original reference}{Original references}
\makeatletter
\newlabel{algebraic-lemma-base-change-representable-by-spaces}{{0302}{source}{}{sourcepointer-2}{}}
\newlabel{algebraic-lemma-base-change-representable-by-spaces@cref}{{[originalreference][0][]0302}{[1][0][]source}{}{}{}}
\makeatother
\crefname{originalreference}{original reference}{original references}
\Crefname{originalreference}{Original reference}{Original references}
\makeatletter
\newlabel{algebraic-lemma-base-change-representable-transformations-property}{{045C}{source}{}{sourcepointer-3}{}}
\newlabel{algebraic-lemma-base-change-representable-transformations-property@cref}{{[originalreference][0][]045C}{[1][0][]source}{}{}{}}
\makeatother
\crefname{originalreference}{original reference}{original references}
\Crefname{originalreference}{Original reference}{Original references}
\makeatletter
\newlabel{algebraic-lemma-map-fibred-setoids-property}{{04TC}{source}{}{sourcepointer-4}{}}
\newlabel{algebraic-lemma-map-fibred-setoids-property@cref}{{[originalreference][0][]04TC}{[1][0][]source}{}{}{}}
\makeatother
\crefname{originalreference}{original reference}{original references}
\Crefname{originalreference}{Original reference}{Original references}
\makeatletter
\newlabel{algebraic-lemma-smooth-surjective-morphism-implies-algebraic}{{05UL}{source}{}{sourcepointer-5}{}}
\newlabel{algebraic-lemma-smooth-surjective-morphism-implies-algebraic@cref}{{[originalreference][0][]05UL}{[1][0][]source}{}{}{}}
\makeatother
\crefname{originalreference}{original reference}{original references}
\Crefname{originalreference}{Original reference}{Original references}
\makeatletter
\newlabel{bootstrap-lemma-representable-by-spaces-over-space}{{02YS}{source}{}{sourcepointer-6}{}}
\newlabel{bootstrap-lemma-representable-by-spaces-over-space@cref}{{[originalreference][0][]02YS}{[1][0][]source}{}{}{}}
\makeatother
\crefname{originalreference}{original reference}{original references}
\Crefname{originalreference}{Original reference}{Original references}
\makeatletter
\newlabel{bootstrap-theorem-final-bootstrap}{{04S6}{source}{}{sourcepointer-7}{}}
\newlabel{bootstrap-theorem-final-bootstrap@cref}{{[originalreference][0][]04S6}{[1][0][]source}{}{}{}}
\makeatother
\crefname{originalreference}{original reference}{original references}
\Crefname{originalreference}{Original reference}{Original references}
\makeatletter
\newlabel{examples-stacks-lemma-faithful-hilbert}{{05XV}{source}{}{sourcepointer-8}{}}
\newlabel{examples-stacks-lemma-faithful-hilbert@cref}{{[originalreference][0][]05XV}{[1][0][]source}{}{}{}}
\makeatother
\crefname{originallemma}{original lemma}{original lemmas}
\Crefname{originallemma}{Original lemma}{Original lemmas}
\makeatletter
\newlabel{lemma-flat-finite-presentation-surjective-diagonal}{{05XW}{source}{}{sourcepointer-9}{}}
\newlabel{lemma-flat-finite-presentation-surjective-diagonal@cref}{{[originallemma][0][]05XW}{[1][0][]source}{}{}{}}
\makeatother
\crefname{originallemma}{original lemma}{original lemmas}
\Crefname{originallemma}{Original lemma}{Original lemmas}
\makeatletter
\newlabel{lemma-hilbert-stack-of-space}{{05YS}{source}{}{sourcepointer-10}{}}
\newlabel{lemma-hilbert-stack-of-space@cref}{{[originallemma][0][]05YS}{[1][0][]source}{}{}{}}
\makeatother
\crefname{originallemma}{original lemma}{original lemmas}
\Crefname{originallemma}{Original lemma}{Original lemmas}
\makeatletter
\newlabel{lemma-relative-hilbert}{{05YH}{source}{}{sourcepointer-11}{}}
\newlabel{lemma-relative-hilbert@cref}{{[originallemma][0][]05YH}{[1][0][]source}{}{}{}}
\makeatother
\crefname{originalreference}{original reference}{original references}
\Crefname{originalreference}{Original reference}{Original references}
\makeatletter
\newlabel{stacks-lemma-2-fibre-product-gives-stack-in-setoids}{{05UI}{source}{}{sourcepointer-12}{}}
\newlabel{stacks-lemma-2-fibre-product-gives-stack-in-setoids@cref}{{[originalreference][0][]05UI}{[1][0][]source}{}{}{}}
\makeatother
\title{The relative finite Hilbert stack over a flat cover is algebraic}
\author{The Stacks Project Authors}
\date{stacks-project:criteria.tex:local-snapshot-20261012; GNU FDL 1.2 or later}
\begin{document}
\maketitle
\noindent\textit{Editorial note.} No omission markers in-span; proof complete as written.\par\medskip
\noindent\textit{Allowed standard prerequisites (editorial).}\begin{itemize}
\item Algebraicity of H\_d(X). lemma-hilbert-stack-of-space (same chapter, named). Provides the smooth cover of H\_d(X).
\item Setoid descent for fibre products. examples-stacks-lemma-faithful-hilbert, stacks-lemma-2-fibre-product-gives-stack-in-setoids (cross-chapter, named). Stack-in-setoids of the 2-fibre product W.
\item Relative representability. lemma-relative-hilbert (same chapter, named). Makes H\_d(X/Y) -> H\_d(X) x Y representable.
\item Bootstrap. bootstrap-theorem-final-bootstrap (cross-chapter, named). Concludes algebraicity of W from its sheaf.
\end{itemize}\medskip
\ifdefined\hypertarget\hypertarget{sourcepointer-1}{}\fi
\noindent\textit{Source reference.} Original reference 0300: focused-seeds/geometry/S9-crit-hilbert-stack-relative-space/ref-tags.txt, source lines 5--5. Navigation pointer only: the counted proof cites this statement by \textbackslash{}ref as a standing input; the deduction the counted proof has to make is carried out inside the counted span, and no part of that deduction is delegated to this pointer.\par\medskip
\ifdefined\hypertarget\hypertarget{sourcepointer-2}{}\fi
\noindent\textit{Source reference.} Original reference 0302: focused-seeds/geometry/S9-crit-hilbert-stack-relative-space/ref-tags.txt, source lines 6--6. Navigation pointer only: the counted proof cites this statement by \textbackslash{}ref as a standing input; the deduction the counted proof has to make is carried out inside the counted span, and no part of that deduction is delegated to this pointer.\par\medskip
\ifdefined\hypertarget\hypertarget{sourcepointer-3}{}\fi
\noindent\textit{Source reference.} Original reference 045C: focused-seeds/geometry/S9-crit-hilbert-stack-relative-space/ref-tags.txt, source lines 7--7. Navigation pointer only: the counted proof cites this statement by \textbackslash{}ref as a standing input; the deduction the counted proof has to make is carried out inside the counted span, and no part of that deduction is delegated to this pointer.\par\medskip
\ifdefined\hypertarget\hypertarget{sourcepointer-4}{}\fi
\noindent\textit{Source reference.} Original reference 04TC: focused-seeds/geometry/S9-crit-hilbert-stack-relative-space/ref-tags.txt, source lines 8--8. Navigation pointer only: the counted proof cites this statement by \textbackslash{}ref as a standing input; the deduction the counted proof has to make is carried out inside the counted span, and no part of that deduction is delegated to this pointer.\par\medskip
\ifdefined\hypertarget\hypertarget{sourcepointer-5}{}\fi
\noindent\textit{Source reference.} Original reference 05UL: focused-seeds/geometry/S9-crit-hilbert-stack-relative-space/ref-tags.txt, source lines 9--9. Navigation pointer only: the counted proof cites this statement by \textbackslash{}ref as a standing input; the deduction the counted proof has to make is carried out inside the counted span, and no part of that deduction is delegated to this pointer.\par\medskip
\ifdefined\hypertarget\hypertarget{sourcepointer-6}{}\fi
\noindent\textit{Source reference.} Original reference 02YS: focused-seeds/geometry/S9-crit-hilbert-stack-relative-space/ref-tags.txt, source lines 10--10. Navigation pointer only: the counted proof cites this statement by \textbackslash{}ref as a standing input; the deduction the counted proof has to make is carried out inside the counted span, and no part of that deduction is delegated to this pointer.\par\medskip
\ifdefined\hypertarget\hypertarget{sourcepointer-7}{}\fi
\noindent\textit{Source reference.} Original reference 04S6: focused-seeds/geometry/S9-crit-hilbert-stack-relative-space/ref-tags.txt, source lines 11--11. Navigation pointer only: the counted proof cites this statement by \textbackslash{}ref as a standing input; the deduction the counted proof has to make is carried out inside the counted span, and no part of that deduction is delegated to this pointer.\par\medskip
\ifdefined\hypertarget\hypertarget{sourcepointer-8}{}\fi
\noindent\textit{Source reference.} Original reference 05XV: focused-seeds/geometry/S9-crit-hilbert-stack-relative-space/ref-tags.txt, source lines 12--12. Navigation pointer only: the counted proof cites this statement by \textbackslash{}ref as a standing input; the deduction the counted proof has to make is carried out inside the counted span, and no part of that deduction is delegated to this pointer.\par\medskip
\ifdefined\hypertarget\hypertarget{sourcepointer-9}{}\fi
\noindent\textit{Source reference.} Original lemma 05XW: sources/stacks-criteria/pinned-original-source.tex, source lines 412--412. Navigation pointer only: the counted proof cites this statement by \textbackslash{}ref as a standing input; the deduction the counted proof has to make is carried out inside the counted span, and no part of that deduction is delegated to this pointer.\par\medskip
\ifdefined\hypertarget\hypertarget{sourcepointer-10}{}\fi
\noindent\textit{Source reference.} Original lemma 05YS: sources/stacks-criteria/pinned-original-source.tex, source lines 2592--2592. Navigation pointer only: the counted proof cites this statement by \textbackslash{}ref as a standing input; the deduction the counted proof has to make is carried out inside the counted span, and no part of that deduction is delegated to this pointer.\par\medskip
\ifdefined\hypertarget\hypertarget{sourcepointer-11}{}\fi
\noindent\textit{Source reference.} Original lemma 05YH: sources/stacks-criteria/pinned-original-source.tex, source lines 2116--2116. Navigation pointer only: the counted proof cites this statement by \textbackslash{}ref as a standing input; the deduction the counted proof has to make is carried out inside the counted span, and no part of that deduction is delegated to this pointer.\par\medskip
\ifdefined\hypertarget\hypertarget{sourcepointer-12}{}\fi
\noindent\textit{Source reference.} Original reference 05UI: focused-seeds/geometry/S9-crit-hilbert-stack-relative-space/ref-tags.txt, source lines 13--13. Navigation pointer only: the counted proof cites this statement by \textbackslash{}ref as a standing input; the deduction the counted proof has to make is carried out inside the counted span, and no part of that deduction is delegated to this pointer.\par\medskip
\begin{lemma}
\label{lemma-hilbert-stack-relative-space}
Let $S$ be a scheme. Let $F : \mathcal{X} \to \mathcal{Y}$ be a $1$-morphism
of stacks in groupoids over $(\Sch/S)_{fppf}$ such that
\begin{enumerate}
\item $\mathcal{X}$ is representable by an algebraic space, and
\item $F$ is representable by algebraic spaces, surjective, flat, and
locally of finite presentation.
\end{enumerate}
Then $\mathcal{H}_d(\mathcal{X}/\mathcal{Y})$ is an algebraic stack.
\end{lemma}

\begin{proof}
Choose a representable stack in groupoids $\mathcal{U}$ over $S$ and a
$1$-morphism $f : \mathcal{U} \to \mathcal{H}_d(\mathcal{X})$
which is representable by algebraic spaces, smooth, and surjective.
This is possible because $\mathcal{H}_d(\mathcal{X})$ is an algebraic stack by
Lemma \ref{lemma-hilbert-stack-of-space}.
Consider the $2$-fibre product
$$
\mathcal{W} =
\mathcal{H}_d(\mathcal{X}/\mathcal{Y})
\times_{\mathcal{H}_d(\mathcal{X}), f}
\mathcal{U}.
$$
Since $\mathcal{U}$ is representable (in particular a stack in setoids)
it follows from
Examples of Stacks, Lemma \ref{examples-stacks-lemma-faithful-hilbert}
and
Stacks, Lemma \ref{stacks-lemma-2-fibre-product-gives-stack-in-setoids}
that $\mathcal{W}$ is a stack in setoids. The $1$-morphism
$\mathcal{W} \to \mathcal{H}_d(\mathcal{X}/\mathcal{Y})$ is
representable by algebraic spaces, smooth, and surjective as a base
change of the morphism $f$ (see
Algebraic Stacks,
Lemmas \ref{algebraic-lemma-base-change-representable-by-spaces} and
\ref{algebraic-lemma-base-change-representable-transformations-property}).
Thus, if we can show that $\mathcal{W}$ is representable by an algebraic space,
then the lemma follows from
Algebraic Stacks,
Lemma \ref{algebraic-lemma-smooth-surjective-morphism-implies-algebraic}.

\medskip\noindent
The diagonal of $\mathcal{Y}$ is representable by algebraic spaces according to
Lemma \ref{lemma-flat-finite-presentation-surjective-diagonal}.
We may apply
Lemma \ref{lemma-relative-hilbert}
to see that the $1$-morphism
$$
\mathcal{H}_d(\mathcal{X}/\mathcal{Y})
\longrightarrow
\mathcal{H}_d(\mathcal{X}) \times \mathcal{Y}
$$
is representable by algebraic spaces. Consider the $2$-fibre product
$$
\mathcal{V} =
\mathcal{H}_d(\mathcal{X}/\mathcal{Y})
\times_{(\mathcal{H}_d(\mathcal{X}) \times \mathcal{Y}), f \times F}
(\mathcal{U} \times \mathcal{X}).
$$
The projection morphism $\mathcal{V} \to \mathcal{U} \times \mathcal{X}$
is representable by algebraic spaces as a base change of the last
displayed morphism. Hence $\mathcal{V}$ is an algebraic space (see
Bootstrap, Lemma \ref{bootstrap-lemma-representable-by-spaces-over-space}
or
Algebraic Stacks,
Lemma \ref{algebraic-lemma-base-change-by-space-representable-by-space}).
The $1$-morphism $\mathcal{V} \to \mathcal{U}$ fits into the following
$2$-cartesian diagram
$$
\xymatrix{
\mathcal{V} \ar[d] \ar[r] & \mathcal{X} \ar[d]^F \\
\mathcal{W} \ar[r] & \mathcal{Y}
}
$$
because
$$
\mathcal{H}_d(\mathcal{X}/\mathcal{Y})
\times_{(\mathcal{H}_d(\mathcal{X}) \times \mathcal{Y}), f \times F}
(\mathcal{U} \times \mathcal{X})
=
(\mathcal{H}_d(\mathcal{X}/\mathcal{Y})
\times_{\mathcal{H}_d(\mathcal{X}), f}
\mathcal{U}) \times_{\mathcal{Y}, F} \mathcal{X}.
$$
Hence $\mathcal{V} \to \mathcal{W}$ is representable by algebraic spaces,
surjective, flat, and locally of finite presentation as a base change
of $F$. It follows that the same thing is true for the corresponding
sheaves of sets associated to $\mathcal{V}$ and $\mathcal{W}$, see
Algebraic Stacks, Lemma \ref{algebraic-lemma-map-fibred-setoids-property}.
Thus we conclude that the sheaf associated to $\mathcal{W}$ is an
algebraic space by
Bootstrap, Theorem \ref{bootstrap-theorem-final-bootstrap}.
\end{proof}

\end{document}
